\documentclass[12pt]{article}

\usepackage{geometry}
\usepackage{titling}
\usepackage{amsmath,amssymb}
\usepackage{fontspec}
\usepackage{unicode-math}
\usepackage{array}
\usepackage{colortbl}
\usepackage[dvipsnames]{xcolor}
\usepackage[most]{tcolorbox}
\usetikzlibrary{calc, angles, quotes}

\geometry{a4paper, top=2.5cm, bottom=2.5cm, left=2.6cm, right=2.5cm}
\setlength{\droptitle}{-3cm}

\setmonofont{DejaVu Sans Mono}
\setmainfont{Libertinus Serif}
\setsansfont{Libertinus Sans}
\setmathfont{STIXTwoMath-Regular.otf}

\newcommand{\MVec}[1]{\mathbf{#1}}
\newcommand{\vsp}{\vspace*{1em}}

\title{The Law of Cosines \vspace*{2pt}\\ and \vspace*{2pt}\\ the Angle Between Two Vectors}
\author{Devi Prasad M}
\date{September 2026}

\begin{document}
\pagecolor{yellow!25!brown!30}

\maketitle


\section{The Law of Cosines}
Consider a triangle with side lengths $a$, $b$, and $c$, and
an angle $\theta$ opposite to side $c$:

\begin{figure}[h]
\centering
\begin{tikzpicture}[scale=1.2]
    \coordinate (A) at (0,0);
    \coordinate (B) at (5,0);
    \coordinate (C) at (2,3);

    % The triangle
    \draw[thick] (A) -- (B) node[midway, below=2pt] {$b$} -- (C) node[midway, above right] {$c$} -- cycle node[midway, above left] {$a$};

    % Angle theta
    \pic [draw, angle radius=0.6cm, "$\theta$", angle eccentricity=1.4] {angle = B--A--C};

    % Vertex labels
    \node[left] at (A) {$A$};
    \node[right] at (B) {$B$};
    \node[above] at (C) {$C$};
\end{tikzpicture}
\caption{A Generic Triangle.}
\label{fig:triangle_diagram}
\end{figure}

\noindent In such a triangle, the following relationship holds:
\begin{equation}
c^2 = a^2 + b^2 - 2ab \cos\theta
\end{equation}

\subsection{The Geometric Construction}

Consider the triangle shown in Figure \ref{fig:triangle_diagram}.
We drop a height (altitude) $h$ perpendicular to base $b$
from the top vertex, dividing base $b$ into two parts of lengths
$x$ and $b - x$.

\begin{figure}[h]
\centering
\begin{tikzpicture}[scale=1.2]
    \coordinate (A) at (0,0);
    \coordinate (B) at (5,0);
    \coordinate (C) at (2,3);
    \coordinate (D) at (2,0);

    \draw[thick] (A) -- (B) node[midway, below=12pt] {$b$} -- (C) node[midway, above right] {$c$} -- cycle node[midway, above left] {$a$};

    \draw[dashed, blue, thick] (C) -- (D) node[midway, right] {$h$};

    % Right angle symbol
    \draw (2,0.25) -- (2.25,0.25) -- (2.25,0);

    \pic [draw, angle radius=0.6cm, "$\theta$", angle eccentricity=1.4] {angle = B--A--C};

    % Segment markings (fixed)
    \draw (A) -- (D) node[midway, below] {$x$};
    \draw (D) -- (B) node[midway, below] {$b-x$};

    \node[left] at (A) {$A$};
    \node[right] at (B) {$B$};
    \node[above] at (C) {$C$};
    \node[below] at (D) {$D$};
\end{tikzpicture}
\caption{The $\triangle \, ABC$ with perpendicular$h$ from vertex $C$ to base $AB$.}
\label{fig:vector_triangle}
\end{figure}

\section{Proof of the Law of Cosines}

\subsection*{Step 1: Express $x$ and $h$ using Trigonometry}
Using the right-angled triangle $\triangle \, ADC$:
\begin{align}
    \cos\theta &= \frac{x}{a} \implies x = a \cos\theta \\
    \sin\theta &= \frac{h}{a} \implies h = a \sin\theta
\end{align}

\subsection*{Step 2: Apply Pythagorean Theorem}
Applying the Pythagorean theorem to the right-angled triangle \, $\triangle BDC$:
\begin{equation}
    c^2 = h^2 + (b - x)^2
\end{equation}

\subsection*{Step 3: Substitute and Simplify}
Substitute $x = a \cos\theta$ and $h = a \sin\theta$ into equation (4):
\begin{align*}
    c^2 &= (a \sin\theta)^2 + (b - a \cos\theta)^2 \\
        &= a^2 \sin^2\theta + \left(b^2 - 2ab \cos\theta + a^2 \cos^2\theta\right) \\
        &= a^2 \left(\sin^2\theta + \cos^2\theta\right) + b^2 - 2ab \cos\theta
\end{align*}

\noindent Using the identity $\sin^2\theta + \cos^2\theta = 1$:
\begin{equation}
    c^2 = a^2 + b^2 - 2ab \cos\theta \quad \blacksquare
\end{equation}

\newpage
\section{From the Law of Cosines to the Inner Product}

\subsection*{Setting up the correspondence}
Let $\MVec{x}$ and $\MVec{y}$ be vectors based at the origin $A$, with
$\MVec{x}$ pointing to $C$ and $\MVec{y}$ pointing to $B$, so that
\[
    a = |\MVec{x}|, \qquad b = |\MVec{y}|,
\]
and $\theta$ is the angle between them, exactly as in
Figure~\ref{fig:triangle_diagram}.

\begin{figure}[h]
\centering
\begin{tikzpicture}[scale=1.2]
    \coordinate (A) at (0,0);
    \coordinate (B) at (5,0);
    \coordinate (C) at (2,3);

    % The triangle, with AB and AC drawn as vectors
    \draw[thick, ->] (A) -- (B) node[midway, below=2pt] {$\MVec{y}$};
    \draw[thick, ->] (A) -- (C) node[midway, above left] {$\MVec{x}$};
    \draw[thick] (B) -- (C) node[midway, above right] {$\MVec{x}-\MVec{y}$};

    % Angle theta
    \pic [draw, angle radius=0.6cm, "$\theta$", angle eccentricity=1.4] {angle = B--A--C};

    % Vertex labels
    \node[left] at (A) {$A$};
    \node[right] at (B) {$B$};
    \node[above] at (C) {$C$};
\end{tikzpicture}
\caption{The vectors $\MVec{x}$ and $\MVec{y}$, with angle $\theta$ between them.}
\label{fig:vector_triangle}
\end{figure}

The third side of the triangle is the
segment $BC$, which as a vector is $\MVec{x} - \MVec{y}$. Its length is
precisely the side $c$:
\begin{equation}
    c = |\MVec{x} - \MVec{y}|.
\end{equation}

\subsection*{Step 4: Expand $c^2$ algebraically}
We use the \emph{algebraic} definition of the inner product,
$\MVec{u}\cdot\MVec{v} = \sum_i u_i v_i$, together with the identity
$|\MVec{u}|^2 = \MVec{u}\cdot\MVec{u}$. Expanding by bilinearity:

\begin{align}
    c^2 = |\MVec{x}-\MVec{y}|^2
        &= (\MVec{x}-\MVec{y})\cdot(\MVec{x}-\MVec{y}) \notag\\
        &= \MVec{x}\cdot\MVec{x} - 2\,\MVec{x}\cdot\MVec{y} + \MVec{y}\cdot\MVec{y} \notag\\
        &= a^2 - 2\,\MVec{x}\cdot\MVec{y} + b^2.
        \label{eq:algebraic-c2}
\end{align}

This is a purely algebraic expression for $c^2$ — no angle appears in it.

\subsection*{Step 5: Equate to the Law of Cosines}
Equation (5) gave us a purely \emph{geometric} expression for the same
quantity $c^2$:

\begin{equation}
    c^2 = a^2 + b^2 - 2ab\cos\theta.
    \label{eq:geometric-c2}
\end{equation}

Since \eqref{eq:algebraic-c2} and \eqref{eq:geometric-c2} both equal $c^2$,
we may set them equal to each other:
\[
    a^2 - 2\,\MVec{x}\cdot\MVec{y} + b^2 = a^2 + b^2 - 2ab\cos\theta.
\]

\subsection*{Step 6: Cancel and solve}
The terms $a^2$ and $b^2$ appear on both sides and cancel:
\[
    -2\,\MVec{x}\cdot\MVec{y} = -2ab\cos\theta.
\]
Dividing both sides by $-2$:

\begin{equation}
    \MVec{x}\cdot\MVec{y} = ab\cos\theta = |\MVec{x}|\,|\MVec{y}|\cos\theta.
    \quad\blacksquare
\end{equation}

\subsection*{Remark}
The Law of Cosines was proved purely from Pythagoras, with no reference to
inner products.

\vsp

\noindent Equation \eqref{eq:algebraic-c2} was proved purely from the
algebraic (coordinate) definition of the inner product, with no reference
to angles.

\vsp

\noindent The bridge between algebra and geometry is the single quantity
$c = |\MVec{x}-\MVec{y}|$: one formula for $c^2$ comes from coordinates,
the other from Pythagoras, and equating them is precisely what encodes
$\cos\theta$ inside the inner product.

\end{document}